\section{总结与延伸}

\subsection{本讲核心要点}

\begin{enumerate}
    \item \textbf{数据鸿沟问题}：图像数据集达数十亿级别，但 3D 模型等其他模态的数据规模小得多。如何利用 2D 先验弥补 3D 数据的不足是核心动机。

    \item \textbf{Neural Rendering 作为桥梁}：可微分渲染函数 $g(\theta, \pi)$ 将 3D 表示映射为 2D 图像，使得 2D 空间的优化信号可以反向传播到 3D 参数。

    \item \textbf{从 CLIP 到扩散模型}：CLIP 提供语义级别的对齐，但细节不足。扩散模型包含更丰富的像素级先验，SDS 提供了利用这些先验的有效方式。

    \item \textbf{SDS 的核心机制}：将预训练扩散模型的噪声预测损失作为渲染质量的度量，通过忽略 U-Net Jacobian 的简化实现高效优化。

    \item \textbf{DreamFusion 的成功与局限}：展示了 SDS 在文本到 3D 生成中的巨大潜力，但 Janus Problem 等问题揭示了 2D 先验的固有局限性。

    \item \textbf{未来趋势}：自回归与扩散模型的边界正在模糊，统一的多模态生成框架是重要方向。
\end{enumerate}

\subsection{方法对照表}

\begin{table}[H]
\centering
\begin{tabular}{p{0.18\textwidth} p{0.22\textwidth} p{0.24\textwidth} p{0.22\textwidth}}
\toprule
方法 & 核心优势 & 主要局限 & 适用场景 \\
\midrule
CLIP 引导 3D & 语义对齐直接、实现相对简单 & 细节弱、容易形状漂移 & 零样本概念探索 \\
SDS / DreamFusion & 能复用强 2D 扩散先验 & Janus problem、过饱和、几何不稳定 & 文本到 3D 的高保真原型生成 \\
VSD / 改进型蒸馏 & 多样性和稳定性更好 & 训练与优化更复杂 & 追求更高质量和更少伪影的 3D 生成 \\
多模态统一生成框架 & 有机会打通 2D/3D/视频边界 & 仍处于快速演化期 & 下一代通用生成模型研究 \\
\bottomrule
\end{tabular}
\caption{Lecture 13 中几类核心方法的比较}
\end{table}

\subsection{进一步思考的问题}

\begin{itemize}
    \item 2D 先验究竟能在多大程度上替代真实 3D 数据，而不会永久受限于视角一致性问题？
    \item 当自回归模型和扩散模型的边界持续模糊时，未来的 3D 生成系统是否还会保留今天清晰的方法分类？
    \item Neural Rendering 在未来更多扮演训练桥梁，还是会成为统一多模态生成中的标准中间表示？
\end{itemize}

\subsection{拓展阅读}

\begin{itemize}
    \item Poole, B. et al. ``DreamFusion: Text-to-3D using 2D Diffusion.'' \textit{ICLR 2023}. \\
    ——首次提出 SDS 用于文本到 3D 生成的开创性工作

    \item Wang, Z. et al. ``ProlificDreamer: High-Fidelity and Diverse Text-to-3D Generation with Variational Score Distillation.'' \textit{NeurIPS 2023}. \\
    ——提出 VSD（Variational Score Distillation），显著改善了 SDS 的过饱和与多样性问题

    \item Mildenhall, B. et al. ``NeRF: Representing Scenes as Neural Radiance Fields for View Synthesis.'' \textit{ECCV 2020}. \\
    ——Neural Radiance Fields 的原始论文，Neural Rendering 的里程碑

    \item Kerbl, B. et al. ``3D Gaussian Splatting for Real-Time Radiance Field Rendering.'' \textit{SIGGRAPH 2023}. \\
    ——3D Gaussian Splatting，一种高效的可微分渲染方法

    \item Jain, A. et al. ``DreamFields: Zero-Shot Text-Guided Object Generation with Dream Fields.'' \textit{CVPR 2022}. \\
    ——CLIP 引导的零样本 3D 生成

    \item Radford, A. et al. ``Learning Transferable Visual Models From Natural Language Supervision.'' \textit{ICML 2021}. \\
    ——CLIP 模型原始论文

    \item Luo, T. et al. ``Discrete Diffusion Modeling by Estimating the Ratios of the Data Distribution.'' \textit{ICML 2024}. \\
    ——离散扩散模型的代表性工作
\end{itemize}

%% ==================== 正文结束 ====================

\end{document}
